\chapter{The Linear Algebra of Machine Learning}

In the previous chapter on numerical linear algebra, we saw what happens when abstract
formulas meet physical machines: numbers get rounded, ill-conditioned systems amplify
noise, and matrices with millions of entries cannot be inverted by exact symbolic
formulas. Real computation demands stable projections, low-rank approximations, and
iterative algorithms that step toward answers one matrix--vector product at a time.

This chapter brings those numerical and geometric tools together in a concrete setting.
Almost every method in modern machine learning -- including recommendation engines,
computer vision, semantic search, and large language models -- is a direct application
of numerical and geometric linear algebra.

When a neural network processes an image, the image is represented as a vector of
pixel intensities and multiplied by weight matrices. When a search engine measures
whether ``king'' and ``queen'' are related, it computes a dot product. When a model
is trained, gradient descent steps downhill on a loss surface along a vector of
partial derivatives.

The aim here is to make the linear algebra under machine learning visible enough that the
field stops looking mysterious and starts looking like \emph{geometry you understand}. We
go slowly, build every idea from a picture, and keep every example small enough to check
by hand.

\begin{focusbox}
The one sentence to carry through the whole chapter: \textbf{a machine learning model
is a function from vectors to vectors, built mostly out of matrix multiplications, and
training is the search for the matrices that make that function fit the data.} Hold onto
that and everything below is a variation on it.
\end{focusbox}

\section{Data is a matrix}
\label{ch34:data-matrix}

Before we can do anything clever, we have to agree on how data even \emph{enters} the
machinery. Everything downstream rests on this, so let us be very concrete.

Suppose you run a small clinic and you have recorded, for each of four patients, their
age, their resting heart rate, and their systolic blood pressure. You could write this
as a little table:

\begin{center}
\begin{tabular}{lccc}
\toprule
patient & age & heart rate & blood pressure \\
\midrule
Ana   & 30 & 60 & 110 \\
Ben   & 50 & 72 & 130 \\
Cleo  & 40 & 68 & 120 \\
Dev   & 60 & 80 & 145 \\
\bottomrule
\end{tabular}
\end{center}

What remains is a $4\times 3$ array of numbers:
\[
\mat{X} =
\begin{bmatrix}
30 & 60 & 110 \\
50 & 72 & 130 \\
40 & 68 & 120 \\
60 & 80 & 145
\end{bmatrix}.
\]
This object -- the \emph{data matrix} -- is where machine learning begins. The
convention almost everyone uses is: \textbf{one row per example, one column per
feature.} Each patient is a point; each measurement is a coordinate.

\begin{intuitionbox}
A data matrix has two complementary readings, and you should be able to flip between
them at will.

\emph{Read it by rows} and each row $\vec{x}_i \in \R^{3}$ is a single patient: a point
floating in a three-dimensional ``feature space'' whose axes are age, heart rate, and
blood pressure. Four patients give you a cloud of four points.

\emph{Read it by columns} and each column is one feature measured across the whole
population: the vector $(30,50,40,60)$ is ``everyone's age.'' Now each column lives in
$\R^{4}$, the ``sample space,'' and you can ask whether two columns point in similar
directions -- which, you will see, is exactly asking whether two features are
correlated.
\end{intuitionbox}

This double reading drives the whole subject. The geometry of the rows
(points in feature space) drives clustering, classification, and nearest-neighbor
search. The geometry of the columns (features in sample space) drives correlation,
covariance, and principal component analysis. \emph{Same matrix, two geometries.}

\begin{figure}[H]
\centering
\begin{tikzpicture}[scale=1.0,>=Stealth]
  % axes
  \draw[->,thick] (0,0) -- (5.6,0) node[right]{feature 1 (age)};
  \draw[->,thick] (0,0) -- (0,4.4) node[above]{feature 2 (heart rate)};
  % four data points (rows)
  \fill[vecblue] (1.2,0.8) circle (2.7pt) node[below right=1pt]{\small Ana};
  \fill[vecblue] (3.6,2.4) circle (2.7pt) node[below right=1pt]{\small Ben};
  \fill[vecblue] (2.4,1.7) circle (2.7pt) node[above left=1pt]{\small Cleo};
  \fill[vecblue] (4.6,3.4) circle (2.7pt) node[above left=1pt]{\small Dev};
  % trend
  \draw[targetgreen,thick,dashed] (0.6,0.3) -- (5.1,3.7);
\end{tikzpicture}
\caption{Reading the data matrix by rows: each patient is a point in feature space (only
two of the three features are drawn). Already you can \emph{see} structure -- the points
trend upward together, the geometric fingerprint of a correlation between age and heart
rate.}
\label{ch34:fig-datacloud}
\end{figure}

\begin{trapbox}
The orientation convention bites everyone at least once. In this book and in most ML
libraries, examples are \emph{rows} and features are \emph{columns}, so $\mat{X}$ is
$n \times d$ ($n$ examples, $d$ features). Classical statistics and Strang's own
least-squares chapter often put the \emph{data vector} $\vec{b}$ as a column and stack
equations as rows -- same math, transposed bookkeeping. Whenever a formula looks
``backwards,'' the first thing to check is whether someone transposed $\mat{X}$. Write
the shapes next to every matrix; shapes never lie.
\end{trapbox}

\subsection{Centering and the mean as a vector}

Almost the first thing you do to a data matrix is \emph{center} it: subtract, from each
column, that column's average. The column means form a vector $\bar{\vec{x}} \in \R^{d}$,
and centering replaces every row $\vec{x}_i$ by $\vec{x}_i - \bar{\vec{x}}$. Geometrically
you are sliding the whole cloud of points so its center of mass sits at the origin.

For our clinic, the column means are $\bar{\vec{x}} = (45,\,70,\,126.25)$, and the
centered matrix is
\[
\mat{X}_c =
\begin{bmatrix}
-15 & -10 & -16.25 \\
\phantom{-}5 & \phantom{-}2 & \phantom{-}3.75 \\
-5 & -2 & -6.25 \\
\phantom{-}15 & \phantom{-}10 & \phantom{-}18.75
\end{bmatrix}.
\]
Centering matters because nearly every ``direction'' question -- which way does the data
vary most? do two features move together? -- only makes sense relative to the center.
Without centering, raw vectors $\vec{x}_i$ point from the arbitrary coordinate
origin $(0,\dots,0)$ rather than measuring deviations from the cloud's mean.

\begin{labbox}
Take three points $(1,2),(3,3),(5,7)$ as the rows of a $3\times 2$ matrix. By hand:
(a) compute the column-mean vector; (b) subtract it from each row to center the data;
(c) check that the columns of your centered matrix each sum to zero. That last check is
your safety net -- a centered column \emph{always} sums to zero, so if it does not, you
made an arithmetic slip.
\end{labbox}

\subsection{The Gram matrix: all the dot products at once}

Matrix algebra now provides an immediate advantage. Suppose you want every pairwise dot
product between rows -- a natural thing, because (as the next sections show) dot products
measure similarity. You do \emph{not} loop over pairs. You write a single matrix product:
\[
\mat{G} = \mat{X}\mat{X}\T, \qquad G_{ij} = \vec{x}_i \cdot \vec{x}_j .
\]
The entry in row $i$, column $j$ of $\mat{X}\mat{X}\T$ is the dot product of row $i$ of
$\mat{X}$ with row $j$ of $\mat{X}$ -- that is just the definition of matrix
multiplication, read carefully. So $\mat{G}$, the \emph{Gram matrix}, packages all
$n^2$ similarities into one symmetric matrix.

Flip the order and you get the other geometry. The $d \times d$ matrix
$\mat{X}\T\mat{X}$ collects all pairwise dot products between \emph{columns} (features),
and -- once the data is centered and divided by $n$ -- it is exactly the
\textbf{covariance matrix} you may have met in statistics. We will return to it; for now
notice the elegant duality: $\mat{X}\mat{X}\T$ compares examples, $\mat{X}\T\mat{X}$
compares features, and they share the same nonzero eigenvalues. The SVD chapter explains
why; for now, notice that one little matrix product, taken two ways, generates both
geometries of your data.

\section{Linear regression: geometric least squares}
\label{ch34:linreg}

The oldest learning algorithm in data analysis is linear regression, and it is pure linear
algebra -- in fact it is the least-squares problem you have already studied, expressed in
the vocabulary of data science. Let us rederive it from scratch from an optimization viewpoint, since the picture carries
straight over to the harder models later.

\subsection{The setup}

We have inputs and we want to predict an output. Say we have the ages and heart rates of
patients and we want to predict blood pressure. We propose the simplest possible model:
the prediction is a \emph{weighted sum} of the inputs, plus an offset:
\[
\hat{y} = w_1 \cdot (\text{age}) + w_2 \cdot (\text{heart rate}) + b.
\]
The numbers $w_1, w_2$ are \emph{weights}; $b$ is the \emph{bias} (statisticians call it
the intercept). Learning means choosing $w_1, w_2, b$ so the predictions $\hat{y}$ match
the true blood pressures as well as possible.

A standard bookkeeping trick absorbs the bias. Append a column of ones to the
data matrix. Then the bias becomes ``just another weight,'' the weight on the constant
feature $1$:
\[
\hat{\vec{y}} = \mat{X}\vec{w}, \qquad
\mat{X} =
\begin{bmatrix}
30 & 60 & 1 \\
50 & 72 & 1 \\
40 & 68 & 1 \\
60 & 80 & 1
\end{bmatrix},
\quad
\vec{w} = \begin{bmatrix} w_1 \\ w_2 \\ b \end{bmatrix}.
\]
Now the entire model is one matrix--vector product, $\hat{\vec{y}} = \mat{X}\vec{w}$. The
vector of all four predictions is $\mat{X}$ acting on the weight vector. \emph{The model
is a linear map from weight space to prediction space.}

\begin{focusbox}
A linear regression model \emph{is} the matrix--vector product $\mat{X}\vec{w}$. Training
it means: given the targets $\vec{y}$, find the $\vec{w}$ that makes $\mat{X}\vec{w}$ as
close to $\vec{y}$ as possible. That is the least-squares problem, and you already know
its answer is a \emph{projection}.
\end{focusbox}

\subsection{Why projection? The column-space picture}

The geometry here is the heart of the whole section. As $\vec{w}$ ranges over
all possible weight vectors, $\hat{\vec{y}} = \mat{X}\vec{w}$ ranges over all linear
combinations of the columns of $\mat{X}$ -- that is, over the \emph{column space}
$\col(\mat{X})$. This is a flat subspace (a plane, here a 3-dimensional one) sitting
inside $\R^{4}$, the space where the target $\vec{y}$ lives.

The target $\vec{y}$ -- the actual blood pressures -- almost never lies exactly in that
subspace. (If it did, the model would fit perfectly, and real data never cooperates.)
So the best we can do is find the point \emph{in} the column space that is closest to
$\vec{y}$. The closest point in a subspace to an outside point is the \textbf{orthogonal
projection}. The error $\vec{y} - \hat{\vec{y}}$ must be perpendicular to the subspace --
because if it had any component lying \emph{along} the subspace, you could slide
$\hat{\vec{y}}$ in that direction and get closer.

\begin{figure}[H]
\centering
\begin{tikzpicture}[scale=1.0,>=Stealth]
  % the column-space plane
  \fill[sourceblue!12] (-0.3,-0.3) -- (5.2,0.4) -- (6.4,2.4) -- (0.9,1.7) -- cycle;
  \node[sourceblue!70!black] at (5.7,2.0) {\small $\col(\mat{X})$};
  % projection foot
  \coordinate (O) at (1.4,0.7);
  \coordinate (P) at (3.6,1.2);   % yhat in plane
  \coordinate (Y) at (4.1,3.4);   % y above plane
  % y vector
  \draw[->,very thick,vecred] (O) -- (Y) node[above]{\small $\vec{y}$ (target)};
  % yhat
  \draw[->,very thick,vecblue] (O) -- (P) node[below right]{\small $\hat{\vec{y}}=\mat{X}\vec{w}$};
  % residual
  \draw[->,thick,quotientorange] (P) -- (Y) node[midway,right]{\small $\vec{y}-\hat{\vec{y}}$};
  % right angle
  \draw (P) ++(-0.28,0.05) -- ++(0.1,0.26) -- ++(0.26,-0.1);
\end{tikzpicture}
\caption{Least squares as projection. The predictions $\mat{X}\vec{w}$ can only reach
points in the column space (the shaded plane). The best prediction is the foot of the
perpendicular from the target $\vec{y}$. The residual $\vec{y}-\hat{\vec{y}}$ is
orthogonal to the column space -- that orthogonality \emph{is} the normal equations.}
\label{ch34:fig-projection}
\end{figure}

``The residual is perpendicular to every column of $\mat{X}$'' translates directly into
algebra. Perpendicular to every column means $\mat{X}\T(\vec{y}-\mat{X}\vec{w}) = \vec{0}$,
which rearranges to the celebrated \textbf{normal equations}:
\[
\boxed{\;\mat{X}\T\mat{X}\,\vec{w} = \mat{X}\T\vec{y}\;}
\]
and, when $\mat{X}\T\mat{X}$ is invertible, the closed-form solution
\[
\vec{w} = (\mat{X}\T\mat{X})\inv \mat{X}\T \vec{y}.
\]

\begin{formalbox}
\textbf{Least-squares regression.} Given $\mat{X}\in\R^{n\times d}$ with linearly
independent columns and $\vec{y}\in\R^{n}$, the weight vector minimizing the squared
error $\norm{\vec{y}-\mat{X}\vec{w}}^2$ is the unique solution of
$\mat{X}\T\mat{X}\,\vec{w}=\mat{X}\T\vec{y}$, namely
$\vec{w}=(\mat{X}\T\mat{X})\inv\mat{X}\T\vec{y}$. The fitted values are
$\hat{\vec{y}} = \mat{X}\vec{w} = \mat{X}(\mat{X}\T\mat{X})\inv\mat{X}\T\,\vec{y} =
\mat{P}\,\vec{y}$, where $\mat{P}$ is the orthogonal projector onto $\col(\mat{X})$.
\end{formalbox}

\begin{recoverybox}
\textbf{Pause and think.} Suppose the target vector $\vec{y}$ already lies entirely inside
the column space $\col(\mat{X})$. What is the fitted prediction $\hat{\vec{y}}$, what is the
residual vector $\vec{y}-\hat{\vec{y}}$, and what is the minimum squared error? What if $\vec{y}$
is instead strictly orthogonal to $\col(\mat{X})$?
\end{recoverybox}

\subsection{A fully worked tiny fit}

Let us actually solve one, small enough to do on paper. Fit a line $\hat{y}=w x + b$ to
three points: $(1,1),(2,2),(3,2)$. With the ones-column trick,
\[
\mat{X}=\begin{bmatrix} 1 & 1 \\ 2 & 1 \\ 3 & 1 \end{bmatrix},
\qquad
\vec{y}=\begin{bmatrix} 1 \\ 2 \\ 2 \end{bmatrix},
\qquad
\vec{w}=\begin{bmatrix} w \\ b \end{bmatrix}.
\]
Build the two pieces of the normal equations:
\[
\mat{X}\T\mat{X}=\begin{bmatrix} 1&2&3 \\ 1&1&1 \end{bmatrix}
\begin{bmatrix} 1&1 \\ 2&1 \\ 3&1 \end{bmatrix}
=\begin{bmatrix} 14 & 6 \\ 6 & 3 \end{bmatrix},
\qquad
\mat{X}\T\vec{y}=\begin{bmatrix} 1&2&3 \\ 1&1&1 \end{bmatrix}
\begin{bmatrix} 1\\2\\2 \end{bmatrix}
=\begin{bmatrix} 11 \\ 5 \end{bmatrix}.
\]
Solve $\begin{bmatrix} 14&6\\6&3 \end{bmatrix}\vec{w}=\begin{bmatrix}11\\5\end{bmatrix}$.
The determinant is $14\cdot 3 - 6\cdot 6 = 42-36 = 6$, so
\[
\vec{w}=\frac{1}{6}\begin{bmatrix} 3 & -6 \\ -6 & 14 \end{bmatrix}
\begin{bmatrix} 11 \\ 5 \end{bmatrix}
=\frac{1}{6}\begin{bmatrix} 33-30 \\ -66+70 \end{bmatrix}
=\frac{1}{6}\begin{bmatrix} 3 \\ 4 \end{bmatrix}
=\begin{bmatrix} 1/2 \\ 2/3 \end{bmatrix}.
\]
So the best-fit line is $\hat{y} = \tfrac12 x + \tfrac23$. Now do the self-check the
projection picture promises -- compute the residuals and confirm they are orthogonal to
the column space. Predict at each $x$, being careful with the fractions:
\[
\hat y(1)=\tfrac12+\tfrac23=\tfrac76,\qquad
\hat y(2)=1+\tfrac23=\tfrac53,\qquad
\hat y(3)=\tfrac32+\tfrac23=\tfrac{13}{6}.
\]
The residuals $\vec{y}-\hat{\vec{y}}$ are then
\[
1-\tfrac76=-\tfrac16,\qquad
2-\tfrac53=\tfrac13,\qquad
2-\tfrac{13}{6}=-\tfrac16,
\]
so the residual vector is $\vec{r}=\big(-\tfrac16,\,\tfrac13,\,-\tfrac16\big)$. Two
orthogonality checks, one per column of $\mat{X}$. Dot with the ones-column:
$-\tfrac16+\tfrac13-\tfrac16=0$. Dot with the $x$-column $(1,2,3)$:
$-\tfrac16+\tfrac23-\tfrac12 = -\tfrac16+\tfrac46-\tfrac36 = 0$. Both zero, exactly as the
projection picture promised: the residual is perpendicular to the whole column space. Get
into the habit of running this check -- a single nonzero dot product means an arithmetic
slip is hiding somewhere, and the ``residuals sum to zero'' version (the ones-column
check) is the fastest tripwire you have.

\begin{trapbox}
The residual being orthogonal to the column space is a \emph{theorem}, not a hope, so it
is the perfect self-check. Whenever you finish a least-squares fit, dot the residual with
each column of $\mat{X}$; every dot product must be (numerically) zero. The
constant-feature version of this is the familiar fact that least-squares residuals always
sum to zero when a bias term is included.
\end{trapbox}

\begin{strangbox}
This is Strang's least-squares chapter, retold for data. Strang draws the same column
space, the same perpendicular, the same normal equations $\mat{A}\T\mat{A}\hat{\vec{x}}
=\mat{A}\T\vec{b}$ -- he just calls the unknown $\hat{\vec{x}}$ and the target $\vec{b}$.
\emph{Before} Strang, let this convince you that ``fitting a model'' is nothing more
exotic than dropping a perpendicular onto a subspace. \emph{After} Strang, notice that
the projector $\mat{P}=\mat{X}(\mat{X}\T\mat{X})\inv\mat{X}\T$ is exactly his hat matrix,
and that regularization (next box) is what you reach for when $\mat{X}\T\mat{X}$ is too
close to singular to invert safely.
\end{strangbox}

\subsection{When the columns nearly collide: regularization}

What if two features are almost the same -- say, temperature in Celsius and temperature
in Fahrenheit? Then two columns of $\mat{X}$ are nearly parallel, $\mat{X}\T\mat{X}$ is
nearly singular (a tiny eigenvalue), and $(\mat{X}\T\mat{X})\inv$ blows up: the weights
become enormous and wildly sensitive to noise. The geometric reading is that the column
space is nearly degenerate, so the projection is well-defined but the \emph{coordinates}
of the projection (the weights) are not pinned down.

The usual fix is \textbf{ridge regression}: add a small multiple of the identity
before inverting,
\[
\vec{w} = (\mat{X}\T\mat{X} + \lambda \mat{I})\inv \mat{X}\T\vec{y}, \qquad \lambda > 0.
\]
Adding $\lambda\mat{I}$ lifts every eigenvalue of $\mat{X}\T\mat{X}$ by $\lambda$, so the
smallest eigenvalue can no longer be near zero -- the matrix is comfortably invertible
and the weights stay bounded. You will recognize this as the same trick that stabilizes
any ill-conditioned linear system; in ML it doubles as a way to express a \emph{prior}
that weights should be small.

\section{Vectors that mean something: embeddings and similarity}
\label{ch34:embeddings}

Linear regression assumed the data already arrived as numbers. But much of what we care
about is words, images, users, and products -- things that are not numbers at all. The
central move of modern ML is to \emph{turn those things into vectors}, called
\textbf{embeddings}, in such a way that geometry encodes meaning. Once a thing is a
vector, the entire toolkit of this book applies to it.

\subsection{From words to vectors}

Imagine assigning to every word a point in $\R^{d}$ (in practice $d$ might be 300 or
4096). We want the assignment to have a property: \textbf{words used in similar contexts
land near each other.} ``King'' and ``queen'' should be close; ``king'' and ``yogurt''
should be far. We will not derive the training procedure here -- it is itself a learning
problem -- but we can completely understand how meaning, once embedded, is \emph{read
off} with linear algebra.

The tool is the dot product, normalized into the \textbf{cosine similarity}:
\[
\cos\theta = \frac{\inner{\vec{u}}{\vec{v}}}{\norm{\vec{u}}\,\norm{\vec{v}}}.
\]
This is the cosine of the angle between the two vectors. It is $1$ when they point the
same way, $0$ when perpendicular, $-1$ when opposite. It ignores length entirely and
cares only about \emph{direction} -- which is exactly right for meaning, since a word's
importance (often reflected in vector length) should not change what it means.

\begin{intuitionbox}
Why \emph{cosine} and not plain distance? Because direction is the carrier of meaning.
Two documents, one short and one long, can be ``about the same thing'' even though their
raw vectors have very different lengths. Cosine similarity normalizes the lengths away
and compares only the directions they point -- the topic, not the volume. Geometrically:
project both vectors onto the unit sphere, then measure the angle between them.
\end{intuitionbox}

\begin{recoverybox}
\textbf{Pause and think.} What is the cosine similarity between any nonzero vector
$\vec{v}$ and its negative $-\vec{v}$? What about between $\vec{v}$ and $10\vec{v}$?
Notice how cosine similarity is completely blind to scale, yet acutely sensitive to
whether two vectors point in opposite directions.
\end{recoverybox}

\subsection{A worked similarity}

Take three tiny 2-dimensional ``embeddings'':
\[
\vec{u}_{\text{king}}=\colvec{2\\1},\quad
\vec{u}_{\text{queen}}=\colvec{3\\2},\quad
\vec{u}_{\text{yogurt}}=\colvec{-1\\3}.
\]
Compare king with queen:
\[
\cos\theta = \frac{(2)(3)+(1)(2)}{\sqrt{2^2+1^2}\,\sqrt{3^2+2^2}}
= \frac{8}{\sqrt{5}\,\sqrt{13}} = \frac{8}{\sqrt{65}} \approx \frac{8}{8.06} \approx 0.992.
\]
Nearly $1$: very similar, the two arrows almost overlap. Now king with yogurt:
\[
\cos\theta = \frac{(2)(-1)+(1)(3)}{\sqrt{5}\,\sqrt{10}}
= \frac{1}{\sqrt{50}} \approx \frac{1}{7.07}\approx 0.141.
\]
Close to $0$: nearly perpendicular, nearly unrelated. The geometry has correctly sorted
royalty from breakfast.

\begin{figure}[H]
\centering
\begin{tikzpicture}[scale=0.8,>=Stealth]
  \draw[->,thick] (-2.0,0) -- (4,0) node[right]{};
  \draw[->,thick] (0,-0.3) -- (0,3.8) node[above]{};
  % vectors drawn to match the embeddings exactly: king (2,1), queen (3,2), yogurt (-1,3)
  \draw[->,very thick,vecblue]  (0,0) -- (2,1) node[right]{\small king $(2,1)$};
  \draw[->,very thick,vecblue]  (0,0) -- (3,2) node[above right]{\small queen $(3,2)$};
  \draw[->,very thick,vecred]   (0,0) -- (-1,3) node[above left]{\small yogurt $(-1,3)$};
  % angle between king and yogurt (~82 degrees)
  \draw[quotientorange] (0.62,0.31) arc (26.6:108.4:0.69);
  \node[quotientorange] at (0.35,1.15) {\small $\approx 82^\circ$};
\end{tikzpicture}
\caption{Embeddings make meaning geometric, drawn with the exact vectors above. ``King''
and ``queen'' point almost the same way (cosine $\approx 0.99$, an angle of about
$7^\circ$); ``yogurt'' points off at nearly a right angle (cosine $\approx 0.14$, about
$82^\circ$). Similarity is an angle.}
\label{ch34:fig-cosine}
\end{figure}

\subsection{Word arithmetic and the meaning of differences}

The single most famous demonstration that embeddings capture structure is the analogy
\[
\vec{u}_{\text{king}} - \vec{u}_{\text{man}} + \vec{u}_{\text{woman}} \approx
\vec{u}_{\text{queen}}.
\]
Read it geometrically. The \emph{difference vector} $\vec{u}_{\text{king}} -
\vec{u}_{\text{man}}$ is the displacement that, in this learned space, points along the
``royalty'' axis. Add it to ``woman'' and you arrive near ``queen'' (Figure~\ref{ch34:fig-wordarith}).
That a consistent
direction in the embedding space corresponds to a consistent semantic relationship
(gender, tense, plurality, capital-of-country) is the headline result of word
embeddings -- and it relies entirely on the vector addition and parallelogram rule you met in
the very first chapter.

\begin{figure}[H]
\centering
\begin{tikzpicture}[scale=1.3,>=Stealth]
  % four word points forming a parallelogram
  \coordinate (man)   at (1.0,1.0);
  \coordinate (woman) at (1.0,2.5);
  \coordinate (king)  at (3.0,1.4);
  \coordinate (queen) at (3.0,2.9);
  % the two parallel "royalty" displacements (king-man) and (queen-woman)
  \draw[->,very thick,vecblue] (man)   -- (king)  node[midway,below]{\small $+\,(\text{king}-\text{man})$};
  \draw[->,very thick,vecblue] (woman) -- (queen) node[midway,above]{\small same displacement};
  % the two parallel "gender" displacements (woman-man) and (queen-king)
  \draw[->,thick,quotientorange] (man)  -- (woman) node[midway,left]{\small $+\,(\text{woman}-\text{man})$};
  \draw[->,thick,quotientorange] (king) -- (queen) node[midway,right]{\small same gender shift};
  % points
  \foreach \p/\lab/\pos in {man/man/below left,woman/woman/above left,king/king/below right,queen/queen/above right}{
    \fill[black] (\p) circle (1.6pt) node[\pos=1pt]{\small \lab};}
\end{tikzpicture}
\caption{Word arithmetic is the parallelogram rule. The displacement
$\vec{u}_{\text{king}}-\vec{u}_{\text{man}}$ (blue) is a fixed ``royalty'' direction; the
same blue displacement carries ``woman'' to ``queen.'' Equivalently, the ``gender''
displacement (orange) is shared by the man$\to$woman and king$\to$queen edges. The four
words sit at the corners of a parallelogram, so
$\vec{u}_{\text{king}}-\vec{u}_{\text{man}}+\vec{u}_{\text{woman}}\approx\vec{u}_{\text{queen}}$.}
\label{ch34:fig-wordarith}
\end{figure}

\begin{applicationbox}
\textbf{Semantic search and retrieval.} When you type a question into a modern search
box, or when a retrieval-augmented language model fetches relevant documents, here is the
machinery: the query is embedded into a vector $\vec{q}$, every document is pre-embedded
into a vector $\vec{d}_i$, and the system returns the documents whose embeddings have the
largest cosine similarity to $\vec{q}$. Stack all document vectors as rows of a matrix
$\mat{D}$ and the entire search is one matrix--vector product $\mat{D}\vec{q}$ followed by
``take the largest entries.'' Billions of documents, searched by meaning, reduce to a
matrix hitting a vector. Specialized ``vector databases'' exist for exactly this and
nothing more.
\end{applicationbox}

\begin{connectionbox}
\textbf{Inner-product geometry.} Cosine similarity is the concrete realization of the abstract \emph{inner product}. The Cauchy--Schwarz inequality
$\abs{\inner{\vec{u}}{\vec{v}}} \le \norm{\vec{u}}\norm{\vec{v}}$ is precisely the
statement that cosine lives in $[-1,1]$ -- the angle is always real. And the
``orthogonal $=$ unrelated'' intuition is the same one that makes Fourier analysis work:
there, perpendicular sine waves carry independent information, exactly as perpendicular
embedding directions carry independent meanings. The same inner-product geometry applies across both settings.
\end{connectionbox}

\section{A neural network layer is $\mat{W}\vec{x}+\vec{b}$, then a bend}
\label{ch34:nnlayer}

Now we build the object at the center of modern AI. A neural network is a stack of
\emph{layers}, and a single layer is almost entirely linear algebra. Understand one
layer and you understand the whole stack, because a deep network is just layers fed into
layers.

\subsection{The two-step recipe}

A layer takes an input vector $\vec{x}\in\R^{d}$ and produces an output vector
$\vec{a}\in\R^{m}$ in two steps:
\begin{enumerate}
  \item \textbf{Affine step (linear algebra):} compute $\vec{z} = \mat{W}\vec{x} +
  \vec{b}$, where $\mat{W}\in\R^{m\times d}$ is a weight matrix and $\vec{b}\in\R^{m}$ a
  bias vector.
  \item \textbf{Nonlinear step (the bend):} apply a fixed nonlinear function $\sigma$ to
  each entry, $\vec{a} = \sigma(\vec{z})$.
\end{enumerate}
That is the entire layer: $\vec{a} = \sigma(\mat{W}\vec{x}+\vec{b})$. The matrix $\mat{W}$
does all the mixing; the function $\sigma$ does all the bending.

You understand the affine step deeply by now. $\mat{W}\vec{x}$ takes a
linear combination of the input features (each output is a weighted vote over the inputs),
and $+\vec{b}$ shifts the result. Each row of $\mat{W}$ is one ``detector'': it dots the
input against a stored pattern, large when the input resembles that pattern.

\subsection{Why the bend is non-negotiable}

The most important structural fact about neural networks is a pure linear algebra
argument. \textbf{Without the nonlinearity, a deep network collapses to a single
layer.} Suppose you stacked two affine steps with no bend in between:
\[
\vec{z}_2 = \mat{W}_2(\mat{W}_1 \vec{x} + \vec{b}_1) + \vec{b}_2
= (\mat{W}_2\mat{W}_1)\vec{x} + (\mat{W}_2\vec{b}_1 + \vec{b}_2)
= \mat{W}'\vec{x} + \vec{b}'.
\]
The composition of two affine maps is \emph{again one affine map}, with
$\mat{W}'=\mat{W}_2\mat{W}_1$. Ten layers, a hundred layers -- without bends they all
multiply out into a single matrix. You would have spent enormous effort to build a glorified
linear regression. The nonlinearity $\sigma$ is precisely what breaks this collapse and
lets depth buy you genuinely new, curved decision boundaries.

\begin{focusbox}
\textbf{Matrices mix; nonlinearities bend.} Stacking matrices alone gets you one bigger
matrix (composition of linear maps is linear). The nonlinear $\sigma$ between layers is
the only reason a deep network can do something a single matrix cannot. Linear algebra
supplies the power; the bend supplies the \emph{novelty}.
\end{focusbox}

\begin{historybox}
In 1958, Cornell psychologist Frank Rosenblatt unveiled the Perceptron at the Cornell Aeronautical Laboratory. Built as a five-ton custom computer with motorized potentiometers (variable resistors turned by electric motors to adjust weights), the machine was hailed by the press as the dawn of artificial intelligence. Mathematically, it was a single flat cut through feature space: an affine hyperplane $\vec{w}\T\vec{x} + b \gtrless 0$ separating positive from negative examples.

Eleven years later, MIT mathematicians Marvin Minsky and Seymour Papert published a devastating geometric critique in their book \emph{Perceptrons} (1969). They proved that a single straight boundary could not even compute the Exclusive-OR ($\mathrm{XOR}$) logic gate---a four-point toy problem where opposite corners of a square share the same label. Because no single straight line can separate two diagonally opposite pairs from each other, Minsky and Papert argued that perceptron-like machines were fundamentally limited. A bitter intellectual feud broke out between connectionists, who believed in learning from data, and the symbolic AI establishment, who favored formal logic rules. The symbolic camp prevailed; funding evaporated, and neural research froze for nearly two decades in the First AI Winter.

The thaw came in 1986, when David Rumelhart, Geoffrey Hinton, and Ronald Williams demonstrated that adding hidden layers with smooth, nonlinear bends could solve XOR and discover internal representations. In 1989, Yann LeCun showed that multi-layer networks trained this way could reliably read handwritten digits on postal envelopes. The engine driving this was \emph{backpropagation}: the multivariable chain rule written in matrix notation, where forward data multiplies by weight matrices $\mat{W}$ and backward error gradients multiply by their transposes $\mat{W}\T$.

By 2017, sequential recurrent models---which processed text one slow token at a time, struggling with memory bottlenecks and fading gradients---gave way to the Transformer of Ashish Vaswani and coauthors. Their key move was replacing sequential step-by-step processing with parallel linear projections into query, key, and value subspaces ($\mat{Q} = \mat{X}\mat{W}_Q$, $\mat{K} = \mat{X}\mat{W}_K$, $\mat{V} = \mat{X}\mat{W}_V$) and scaled dot-product attention. When models subsequently exploded to hundreds of billions of parameters, Edward Hu and colleagues introduced Low-Rank Adaptation (LoRA, 2021), freezing the giant weight matrix and learning only a low-rank update $\Delta\mat{W} = \mat{B}\mat{A}$ with $r \ll d$---compressing millions of gradient calculations onto consumer hardware through the Eckart--Young principle.

Yet basic questions remain open. Why gradient descent on non-convex surfaces with far
more parameters than data points reliably discovers flat, generalizable minima rather
than memorizing noise remains an unsolved problem. Equally open is the geometry of
\emph{superposition}---the phenomenon where networks compress more distinct concepts into
activation vectors than there are dimensions by exploiting nearly orthogonal directions
in high dimensions.
\end{historybox}

The most common bend today is the \textbf{ReLU}, ``rectified linear unit,''
$\sigma(z)=\max(0,z)$: it passes positive values through and zeros out negatives. It is
about as gentle a bend as exists -- piecewise linear, just a single kink at the origin --
yet stacking many of them lets the network carve input space into many linear pieces and
fit highly complex functions.

\begin{figure}[H]
\centering
\begin{tikzpicture}[>=Stealth]
\begin{axis}[
  width=11cm, height=5cm,
  axis lines=middle, xlabel={$z$}, ylabel={$\sigma(z)$},
  xmin=-4, xmax=4, ymin=-0.4, ymax=3.2,
  xtick={-4,-2,2,4}, ytick={1,2,3},
  every axis x label/.style={at={(ticklabel* cs:1)},anchor=west},
  every axis y label/.style={at={(ticklabel* cs:1)},anchor=south},
  legend pos=north west, legend cell align=left, samples=100,
]
  \addplot[vecblue,very thick,domain=-4:4] {max(0,x)};
  \addlegendentry{ReLU $\max(0,z)$}
  \addplot[vecred,very thick,domain=-4:4] {1/(1+exp(-x))};
  \addlegendentry{sigmoid $1/(1+e^{-z})$}
\end{axis}
\end{tikzpicture}
\caption{The two most common bends. ReLU (blue) is piecewise linear with a single kink at
the origin: it zeros out negatives and passes positives through unchanged. The sigmoid
(red) smoothly squashes the whole line into $(0,1)$, which is why it reads a linear score
as a probability. The \emph{kink} of ReLU and the \emph{curve} of the sigmoid are the
``novelty'' that pure matrix multiplication cannot supply.}
\label{ch34:fig-activations}
\end{figure}

\subsection{A layer computed by hand}

Let the input be $\vec{x}=\colvec{1\\2}$ and the layer have
\[
\mat{W}=\begin{bmatrix} 1 & -1 \\ 0 & 2 \\ 3 & 1 \end{bmatrix},
\qquad
\vec{b}=\begin{bmatrix} 0 \\ -5 \\ 1 \end{bmatrix},
\qquad \sigma=\text{ReLU}.
\]
First the affine step:
\[
\mat{W}\vec{x}=\begin{bmatrix} 1\cdot1+(-1)\cdot2 \\ 0\cdot1+2\cdot2 \\ 3\cdot1+1\cdot2 \end{bmatrix}
=\begin{bmatrix} -1 \\ 4 \\ 5 \end{bmatrix},
\qquad
\vec{z}=\mat{W}\vec{x}+\vec{b}=\begin{bmatrix} -1 \\ -1 \\ 6 \end{bmatrix}.
\]
Then the bend, applied entrywise:
\[
\vec{a}=\text{ReLU}(\vec{z})=\begin{bmatrix} \max(0,-1) \\ \max(0,-1) \\ \max(0,6) \end{bmatrix}
=\begin{bmatrix} 0 \\ 0 \\ 6 \end{bmatrix}.
\]
Two of the three neurons ``switched off'' (their pre-activations were negative); one
fired with value $6$. That pattern of on/off across neurons is how a network represents
``which features are present'' in this particular input.

\begin{figure}[H]
\centering
\begin{tikzpicture}[scale=0.9,>=Stealth]
  % input nodes
  \foreach \i/\y/\name in {1/2.4/$x_1$,2/0.8/$x_2$}{
    \node[circle,draw,fill=sourceblue!15,minimum size=20pt] (I\i) at (0,\y) {\small \name};}
  % output nodes
  \foreach \j/\y/\name in {1/3.2/$a_1$,2/1.6/$a_2$,3/0/$a_3$}{
    \node[circle,draw,fill=targetgreen!15,minimum size=20pt] (O\j) at (3.4,\y) {\small \name};}
  % edges
  \foreach \i in {1,2}{
    \foreach \j in {1,2,3}{
      \draw[gray!60,->] (I\i) -- (O\j);}}
  \node at (1.7,4.1) {\small $\vec{a}=\sigma(\mat{W}\vec{x}+\vec{b})$};
  \node[align=center] at (1.7,-1.0) {\small each arrow is one weight $W_{ij}$;\\[-2pt] \small each green node adds its bias then bends};
\end{tikzpicture}
\caption{A single fully connected layer. Every input feeds every output through a weight;
the weight matrix $\mat{W}$ \emph{is} this complete bipartite wiring. The affine step
$\mat{W}\vec{x}+\vec{b}$ is the wiring; the nonlinearity is applied inside each green
node.}
\label{ch34:fig-layer}
\end{figure}

\subsection{Classification and logistic regression}

For a yes/no decision (spam or not, tumor or not), we want the model to output a
\emph{probability} in $[0,1]$. The trick is to take a linear score and squash it through
the \textbf{sigmoid} (logistic) function
\[
\sigma(z)=\frac{1}{1+e^{-z}}, \qquad z = \vec{w}\cdot\vec{x}+b.
\]
This is \textbf{logistic regression}, and structurally it is exactly one neuron: an
affine step followed by a bend (the smooth red curve of Figure~\ref{ch34:fig-activations}). The linear part $\vec{w}\cdot\vec{x}+b$ defines a
\emph{hyperplane} in feature space -- the set where the score is zero -- and the sigmoid
turns ``how far are you from the plane, and on which side'' into a probability. Points far
on the positive side get probability near $1$, far on the negative side near $0$, and
right on the plane exactly $0.5$.

\begin{intuitionbox}
A linear classifier is a \emph{ruled line} (in 2D), \emph{plane} (in 3D), or
\emph{hyperplane} (in general) that splits feature space into two half-spaces. The weight
vector $\vec{w}$ is the \textbf{normal} to that hyperplane -- it points in the direction
of ``more positive class.'' The bias $b$ slides the hyperplane toward or away from the
origin. Training a linear classifier is choosing the orientation $\vec{w}$ and the offset
$b$ of one flat cut through the data.
\end{intuitionbox}

\begin{figure}[H]
\centering
\begin{tikzpicture}[scale=0.95,>=Stealth]
  % decision line
  \draw[structurepurple,very thick] (-0.4,2.6) -- (4.4,0.2);
  \node[structurepurple] at (4.6,0.0) {\small $\vec{w}\cdot\vec{x}+b=0$};
  % normal vector
  \draw[->,vecred,very thick] (2,1.4) -- (2.9,2.2) node[above right]{\small $\vec{w}$};
  % positive class
  \fill[targetgreen] (3,2.4) circle (2.5pt);
  \fill[targetgreen] (3.6,2.0) circle (2.5pt);
  \fill[targetgreen] (2.7,3.0) circle (2.5pt);
  \fill[targetgreen] (3.9,2.9) circle (2.5pt);
  \node[targetgreen!60!black] at (4.0,3.4) {\small class $+$};
  % negative class
  \fill[vecblue] (0.6,0.6) circle (2.5pt);
  \fill[vecblue] (1.1,0.3) circle (2.5pt);
  \fill[vecblue] (0.3,1.3) circle (2.5pt);
  \fill[vecblue] (1.5,0.9) circle (2.5pt);
  \node[vecblue!70!black] at (0.2,0.0) {\small class $-$};
\end{tikzpicture}
\caption{A linear classifier. The decision boundary is the hyperplane
$\vec{w}\cdot\vec{x}+b=0$; the weight vector $\vec{w}$ is its normal, pointing toward the
positive class. Logistic regression wraps the signed distance to this plane in a sigmoid
to read it as a probability.}
\label{ch34:fig-classifier}
\end{figure}

\section{Training: gradient descent and the gradient as a vector}
\label{ch34:gd}

We have models -- regression, logistic, layers. We have not yet said how the optimal weights
are determined, except in the one case (linear regression) where a closed-form
projection yields the answer. For general models and loss functions, we rely on the primary optimization method of
modern machine learning: \textbf{gradient descent}. And the gradient is, once more, a vector you can
picture geometrically.

\subsection{The loss surface}

Pick a way to measure how wrong the model is -- a \emph{loss} $L(\vec{w})$, a single
number that depends on all the weights. For regression, the natural loss is the squared
error $L(\vec{w})=\norm{\vec{y}-\mat{X}\vec{w}}^2$. As $\vec{w}$ ranges over weight space,
$L$ traces out a curved surface of hills and valleys sitting above weight space.
Training means \textbf{finding the lowest point on this surface}.

\begin{figure}[H]
\centering
\begin{tikzpicture}[scale=1.0,>=Stealth]
  % a bowl (contours)
  \foreach \r in {0.6,1.2,1.8,2.4}{
    \draw[gray!50] (2,1.3) ellipse ({\r} and {0.62*\r});}
  \node[gray!60!black] at (4.7,2.6) {\small contours of $L$};
  % gradient descent path
  \coordinate (p0) at (4.0,2.6);
  \coordinate (p1) at (3.3,1.95);
  \coordinate (p2) at (2.75,1.6);
  \coordinate (p3) at (2.35,1.42);
  \coordinate (p4) at (2.08,1.34);
  \fill[vecred] (p0) circle (2pt);
  \draw[->,vecred,thick] (p0)--(p1);
  \draw[->,vecred,thick] (p1)--(p2);
  \draw[->,vecred,thick] (p2)--(p3);
  \draw[->,vecred,thick] (p3)--(p4);
  \fill[vecblue] (2,1.3) circle (2.3pt) node[below=2pt]{\small minimum};
  \node[vecred] at (4.5,3.0) {\small start};
  % a downhill arrow = -gradient
  \draw[->,structurepurple,very thick] (3.3,1.95) -- ++(-0.55,-0.35);
  \node[structurepurple] at (2.2,2.25) {\small $-\nabla L$};
\end{tikzpicture}
\caption{Gradient descent. The contours are level sets of the loss $L(\vec{w})$. At each
point the gradient $\nabla L$ points uphill (steepest ascent), so we step in the
direction $-\nabla L$ and roll toward the minimum. Each red arrow is one update
$\vec{w}\leftarrow\vec{w}-\eta\nabla L$.}
\label{ch34:fig-gd}
\end{figure}

\subsection{The gradient points uphill}

The \textbf{gradient} $\nabla L$ is the vector of partial derivatives, one per weight:
\[
\nabla L = \left(\frac{\partial L}{\partial w_1},\,\frac{\partial L}{\partial w_2},\,
\ldots,\,\frac{\partial L}{\partial w_d}\right).
\]
Its defining geometric property is the one to memorize: \textbf{at any point, $\nabla L$
points in the direction of steepest increase of $L$, and its length is the rate of that
increase.} So $-\nabla L$ points in the direction of steepest \emph{decrease} -- straight
downhill. To minimize the loss, repeatedly step a little way downhill:
\[
\boxed{\;\vec{w} \leftarrow \vec{w} - \eta\,\nabla L(\vec{w})\;}
\]
where $\eta>0$ is the \textbf{learning rate}, the step size. That single update rule,
iterated, trains essentially every model in deep learning.

\begin{intuitionbox}
You are standing on a foggy hillside and want to reach the valley. You cannot see far,
but you can feel the slope under your feet. The gradient is the direction the ground rises
fastest; turn around and take a step. Repeat. The learning rate is your stride length:
too small and you crawl; too large and you may overshoot the valley and bounce up the far
slope. Gradient descent turns this local downhill step into an iterative algorithm.
\end{intuitionbox}

\subsection{Why $-\nabla L$ is steepest: a one-line derivation}

This is worth deriving, because it ties the gradient to the dot product. The change in
$L$ when you step a tiny amount $\vec{u}$ (a unit direction) is, to first order, the
\emph{directional derivative}
\[
\frac{d}{dt}\Big|_{t=0} L(\vec{w}+t\vec{u}) = \nabla L \cdot \vec{u}
= \norm{\nabla L}\,\norm{\vec{u}}\cos\theta = \norm{\nabla L}\cos\theta,
\]
using $\norm{\vec{u}}=1$. Among all unit directions $\vec{u}$, this dot product is
\emph{largest} when $\cos\theta=1$, i.e.\ $\vec{u}$ points along $\nabla L$, and
\emph{smallest} (most negative) when $\vec{u}$ points along $-\nabla L$. So $-\nabla L$ is
the steepest-descent direction. The whole justification for gradient descent is one
application of ``the dot product is maximized when the vectors align.''

\subsection{Gradient descent on our regression}

Let us connect the two methods. For $L(\vec{w})=\norm{\vec{y}-\mat{X}\vec{w}}^2$, the
gradient works out to
\[
\nabla L = -2\,\mat{X}\T(\vec{y}-\mat{X}\vec{w}) = 2\big(\mat{X}\T\mat{X}\,\vec{w}-\mat{X}\T\vec{y}\big).
\]
Setting $\nabla L = \vec{0}$ -- the condition for a flat spot, the bottom of the bowl --
gives back exactly the normal equations $\mat{X}\T\mat{X}\vec{w}=\mat{X}\T\vec{y}$! The
closed-form projection answer and the gradient-descent answer are the \emph{same point},
the minimum of the same bowl. Linear regression is the rare case where you can jump to the
bottom in one step; gradient descent is what you use when no such jump exists.

\begin{labbox}
Minimize $L(w)=(w-3)^2$ by hand-running gradient descent. Here $\nabla L = 2(w-3)$. Start
at $w_0=0$ with learning rate $\eta=0.25$. Compute $w_1=w_0-\eta\cdot 2(w_0-3)$, then
$w_2$, then $w_3$. Watch the iterates march toward $3$. Then try $\eta=1.1$ and watch them
\emph{diverge} -- a concrete feel for why the learning rate must not be too large.
\end{labbox}

\section{Backpropagation: the multivariable chain rule as matrix multiplication}
\label{ch34:backprop}

Gradient descent needs the gradient. For a deep network with millions of weights, how do
you compute $\partial L / \partial W_{ij}$ for every single weight, efficiently? The
answer is \textbf{backpropagation}, which is the multivariable chain rule from calculus organized as a
sequence of matrix multiplications by \emph{Jacobians}.

\subsection{The Jacobian: the local linear map}

A neural network is a composition of functions: input $\to$ layer 1 $\to$ layer 2 $\to
\cdots \to$ loss. Each function maps vectors to vectors. The right object to describe how a
vector-to-vector function $\vec{f}:\R^{n}\to\R^{m}$ stretches space \emph{near a point} is
its \textbf{Jacobian matrix} $\mat{J}$, the $m\times n$ matrix of all partial derivatives,
$J_{ij}=\partial f_i/\partial x_j$.

\begin{intuitionbox}
Zoom in far enough on any smooth function and it looks linear -- that linear approximation
is the Jacobian. The Jacobian is the matrix that \emph{is} the derivative of a vector
function: $\vec{f}(\vec{x}+\Delta\vec{x}) \approx \vec{f}(\vec{x}) + \mat{J}\,\Delta\vec{x}$.
Every layer of a network has a Jacobian -- the local matrix that says how a small wiggle
in its input becomes a small wiggle in its output. Notice that for the affine step
$\vec{z}=\mat{W}\vec{x}+\vec{b}$, the Jacobian with respect to $\vec{x}$ is exactly
$\mat{W}$ itself, because an affine map \emph{is} its own linear approximation.
\end{intuitionbox}

\subsection{Chain rule as Jacobian product}

For a composition $\vec{f} = \vec{f}_3 \circ \vec{f}_2 \circ \vec{f}_1$, the multivariable
chain rule says the Jacobian of the whole is the \emph{product} of the Jacobians of the
parts:
\[
\mat{J}_{\vec{f}} = \mat{J}_{3}\,\mat{J}_{2}\,\mat{J}_{1}.
\]
This is the same ``matrices compose by multiplication'' fact that has run through the
entire book, now applied to derivatives. The loss is a scalar at the end of the chain, so
its gradient is obtained by multiplying these Jacobians together. Backpropagation is the
specific, clever \emph{order} in which you do that multiplication.

\subsection{Why multiply starting from the loss (backwards through the net)}

The efficiency insight is purely about matrix-multiplication bookkeeping.
The chain $\mat{J}_{\vec{f}} = \mat{J}_3\mat{J}_2\mat{J}_1$ has the loss as the last
function $\vec{f}_3$, so its Jacobian $\mat{J}_3$ is the \emph{leftmost} factor. Because
the loss is one number, $\mat{J}_3$ is a \emph{row vector} (a $1\times m$ object): the
gradient of a scalar. You can multiply out the product $\mat{J}_3\mat{J}_2\mat{J}_1$
either left-to-right (starting from $\mat{J}_3$) or right-to-left (starting from
$\mat{J}_1$). Multiplying \emph{right-to-left} -- start at the input end and build up a
running product -- forces you to carry around full $m\times n$ matrices at each stage.
Multiplying \textbf{left-to-right, starting from the scalar's row vector} $\mat{J}_3$,
keeps a row vector at every stage:
\[
\underbrace{\mat{J}_3}_{1\times m}\,\underbrace{\mat{J}_2}_{m\times k}
= (\text{row})_{1\times k}, \quad \text{then} \times \mat{J}_1.
\]
A row vector times a matrix is another row vector -- cheap. \textbf{That} is
backpropagation: propagate the gradient \emph{backward} through the network, layer by
layer, as a vector being hit by each layer's Jacobian (in transposed form). The
``backward'' in the name is just the choice to multiply the Jacobian chain in the order
that keeps everything a vector instead of a matrix (Figure~\ref{ch34:fig-backprop}).

\begin{figure}[H]
\centering
\begin{tikzpicture}[scale=0.95,>=Stealth]
  % Forward pass (top track)
  \node[font=\footnotesize\bfseries,text=sourceblue!90!black,anchor=west] at (-0.5,2.4) {Forward pass: data vectors flow forward};

  \node[draw=sourceblue!70,fill=sourceblue!10,rounded corners=3pt,inner sep=4pt,font=\footnotesize,align=center] (x) at (0.8,1.4) {$\vec{x}$ \\ {\scriptsize $d\times 1$}};
  \node[draw=sourceblue!70,fill=sourceblue!10,rounded corners=3pt,inner sep=4pt,font=\footnotesize,align=center] (a1) at (4.0,1.4) {$\vec{a}_1$ \\ {\scriptsize $k\times 1$}};
  \node[draw=sourceblue!70,fill=sourceblue!10,rounded corners=3pt,inner sep=4pt,font=\footnotesize,align=center] (a2) at (7.2,1.4) {$\vec{a}_2$ \\ {\scriptsize $m\times 1$}};
  \node[draw=targetgreen!70,fill=targetgreen!10,rounded corners=3pt,inner sep=4pt,font=\footnotesize,align=center] (loss) at (10.4,1.4) {Loss $L$ \\ {\scriptsize scalar}};

  \draw[->,very thick,sourceblue] (x) -- (a1) node[midway,above,font=\scriptsize]{$\mat{W}_1,\;\sigma$};
  \draw[->,very thick,sourceblue] (a1) -- (a2) node[midway,above,font=\scriptsize]{$\mat{W}_2,\;\sigma$};
  \draw[->,very thick,targetgreen] (a2) -- (loss) node[midway,above,font=\scriptsize]{$L(\vec{a}_2,\vec{y})$};

  % Backward pass (bottom track)
  \node[font=\footnotesize\bfseries,text=structurepurple!90!black,anchor=west] at (-0.5,-0.1) {Backward pass: scalar gradient flows backward as a row vector};

  \node[draw=structurepurple!70,fill=structurepurple!10,rounded corners=3pt,inner sep=4pt,font=\footnotesize,align=center] (gx) at (0.8,-1.1) {$\frac{\partial L}{\partial \vec{x}}$ \\ {\scriptsize $1\times d$}};
  \node[draw=structurepurple!70,fill=structurepurple!10,rounded corners=3pt,inner sep=4pt,font=\footnotesize,align=center] (ga1) at (4.0,-1.1) {$\frac{\partial L}{\partial \vec{a}_1}$ \\ {\scriptsize $1\times k$}};
  \node[draw=structurepurple!70,fill=structurepurple!10,rounded corners=3pt,inner sep=4pt,font=\footnotesize,align=center] (ga2) at (7.2,-1.1) {$\frac{\partial L}{\partial \vec{a}_2}$ \\ {\scriptsize $1\times m$}};
  \node[draw=vecred!70,fill=vecred!10,rounded corners=3pt,inner sep=4pt,font=\footnotesize,align=center] (gloss) at (10.4,-1.1) {$\frac{\partial L}{\partial L} = 1$ \\ {\scriptsize $1\times 1$}};

  \draw[<-,very thick,structurepurple] (gx) -- (ga1) node[midway,below,font=\scriptsize]{$\times\,\mat{J}_1$};
  \draw[<-,very thick,structurepurple] (ga1) -- (ga2) node[midway,below,font=\scriptsize]{$\times\,\mat{J}_2$};
  \draw[<-,very thick,vecred] (ga2) -- (gloss) node[midway,below,font=\scriptsize]{$\times\,\mat{J}_3$};

  % Explanation box at bottom
  \draw[draw=spacegray!30,fill=codebg,rounded corners=3pt] (-0.5,-2.5) rectangle (11.7,-1.8);
  \node[font=\scriptsize,text=spacegray!90!black,align=center] at (5.6,-2.15) {Multiplying left-to-right from the scalar keeps every intermediate gradient a \textbf{row vector},\\avoiding the costly formation of full $m\times d$ intermediate Jacobian matrices.};
\end{tikzpicture}
\caption{Forward evaluation versus backward propagation. Forward: input vector $\vec{x}$ is transformed layer-by-layer into the scalar loss $L$. Backward: the chain rule multiplies Jacobians from the left starting at the scalar output. Each step computes $(\text{row vector})\times(\text{Jacobian matrix}) = (\text{new row vector})$, never materializing a full matrix-matrix product.}
\label{ch34:fig-backprop}
\end{figure}

\begin{focusbox}
\textbf{Forward pass:} push the input vector through the layers, computing each
$\vec{a}=\sigma(\mat{W}\vec{x}+\vec{b})$ and remembering the intermediate values.
\textbf{Backward pass:} push the gradient (a vector) back through the layers, multiplying
by each layer's Jacobian, accumulating $\partial L/\partial \mat{W}$ along the way.
Forward is matrices acting on data; backward is the same matrices (transposed) acting on
gradients. Same linear algebra, run in reverse.
\end{focusbox}

\begin{connectionbox}
\textbf{Backpropagation as the multivariable chain rule.} Everything ``deep learning''
about backprop is engineering -- caching intermediate activations, doing it on a GPU. The
mathematics is the chain rule $\frac{d}{dx}f(g(x))=f'(g(x))g'(x)$ generalized to vectors,
where the single-variable derivatives become Jacobian matrices and ordinary
multiplication becomes matrix multiplication. If you have done a multivariable calculus
course, you have already proven the hard part. Linear algebra supplies the only new idea:
that the chain of derivatives is a chain of matrices, so the \emph{order} of multiplication
is a free choice you can exploit.
\end{connectionbox}

\section{Attention: scaled dot products and matrix weights}
\label{ch34:attention}

The architecture behind modern language models -- the \emph{transformer} --
is built on a mechanism called \textbf{attention}. From a mathematical perspective, attention
is a structure you are now fully equipped to analyze: it is a set of dot products, scaled and
normalized via softmax, used to form a weighted average. Let us assemble it from first principles.

\subsection{Queries, keys, values}

Imagine a sequence of tokens (say, words in a sentence), each already an embedding vector.
Attention lets each token ``look at'' the others and pull in the information it needs. To
do that, every token is turned into three vectors by three learned matrices:
\begin{itemize}
  \item a \textbf{query} $\vec{q}$ -- ``what am I looking for?''
  \item a \textbf{key} $\vec{k}$ -- ``what do I offer?''
  \item a \textbf{value} $\vec{v}$ -- ``what will I hand over if chosen?''
\end{itemize}
Stack the queries as rows of $\mat{Q}$, the keys as rows of $\mat{K}$, the values as rows
of $\mat{V}$. Each is just the embedding matrix times a learned weight matrix -- linear
algebra you have done a dozen times.

\subsection{The core formula, read piece by piece}

The famous formula is
\[
\text{Attention}(\mat{Q},\mat{K},\mat{V}) =
\text{softmax}\!\left(\frac{\mat{Q}\mat{K}\T}{\sqrt{d_k}}\right)\mat{V}.
\]
It looks dense; take it apart and every piece is something from this book.

\textbf{Step 1: the scores $\mat{Q}\mat{K}\T$.} This $n\times n$ matrix is, entry by entry,
$(\mat{Q}\mat{K}\T)_{ij} = \vec{q}_i \cdot \vec{k}_j$ -- the dot product of token $i$'s
query with token $j$'s key. It measures how much token $i$ should care about token $j$.
This is the Gram-matrix idea from Section~\ref{ch34:data-matrix} again: \emph{all pairwise
similarities at once}, as a single matrix product. Large dot product means ``these two are
relevant to each other.''

\textbf{Step 2: the scaling $1/\sqrt{d_k}$.} Dividing by the square root of the key
dimension keeps the dot products from growing too large when the vectors are
high-dimensional (a high-dimensional dot product is a sum of many terms and tends to
spread out). It is a normalization for numerical stability, nothing deeper.

\textbf{Step 3: softmax, row by row.} The softmax turns each row of scores into a
probability distribution -- nonnegative numbers that sum to $1$:
\[
\text{softmax}(\vec{s})_j = \frac{e^{s_j}}{\sum_k e^{s_k}}.
\]
After this, row $i$ holds the \emph{attention weights}: how much of its new representation
token $i$ should draw from each other token. (This is the same exponential-normalization
idea as the sigmoid, generalized to many outputs.)

\textbf{Step 4: the weighted sum $(\cdots)\mat{V}$.} Multiplying the attention-weight
matrix by $\mat{V}$ takes, for each token, a \emph{weighted average of all the value
vectors}, weighted by relevance. A matrix of probabilities times a matrix of values is a
matrix of convex combinations -- output row $i$ is a blend of the value vectors, heaviest
on the tokens token $i$ attended to.

\begin{intuitionbox}
Attention is ``soft lookup.'' A hard lookup (like a dictionary) takes a key and returns
exactly one value. Attention instead compares your query to \emph{every} key by dot
product, turns those similarities into weights that sum to one, and returns a
\emph{blend} of all the values weighted by how well each key matched. The dot product
asks ``how relevant?''; the softmax turns relevance into ``how much weight''; the
multiply by $\mat{V}$ collects the blend. Three operations, all of them linear algebra
plus one exponential.
\end{intuitionbox}

\begin{recoverybox}
\textbf{Pause and think.} Suppose token $1$'s query vector $\vec{q}_1$ happens to be
orthogonal to every key vector $\vec{k}_j$ in the sequence. What are the dot-product
scores $\vec{q}_1\cdot\vec{k}_j$? What attention weights does the softmax produce?
(Because every score is zero, every exponent $e^0$ is $1$, yielding an equal split:
$\alpha_{1j} = 1/n$ across all $n$ tokens.)
\end{recoverybox}

\begin{figure}[H]
\centering
\begin{tikzpicture}[scale=0.88,>=Stealth]
  % Lane 1: Geometric Projections into Semantic Subspaces
  \begin{scope}[shift={(0,4.2)}]
    \node[anchor=south,font=\footnotesize\scshape,text=spacegray!90!black] at (6.6,2.2) {1. Linear Projections into Semantic Subspaces};
    \draw[fill=codebg,draw=spacegray!30,rounded corners=4pt] (-0.6,-1.6) rectangle (14.2,2.1);

    % Token input space
    \node[draw=spacegray!70,fill=white,rounded corners=3pt,inner sep=4pt,font=\footnotesize,align=center] (tok) at (1.2,0.2) {
      \textbf{Token $\vec{x}_i$}\\
      $\vec{x}_1, \vec{x}_2 \in \R^d$
    };

    % Target Subspace Boxes
    \node[draw=sourceblue!70,fill=sourceblue!12,rounded corners=3pt,inner sep=3.5pt,font=\scriptsize,align=center] (qbox) at (5.0,1.35) {
      \textbf{Query Space}\\
      $\vec{q}_i = \mat{W}_Q \vec{x}_i$
    };
    \node[draw=targetgreen!70,fill=targetgreen!12,rounded corners=3pt,inner sep=3.5pt,font=\scriptsize,align=center] (kbox) at (5.0,0.2) {
      \textbf{Key Space}\\
      $\vec{k}_j = \mat{W}_K \vec{x}_j$
    };
    \node[draw=quotientorange!70,fill=quotientorange!12,rounded corners=3pt,inner sep=3.5pt,font=\scriptsize,align=center] (vbox) at (5.0,-0.95) {
      \textbf{Value Space}\\
      $\vec{v}_j = \mat{W}_V \vec{x}_j$
    };

    % Projection arrows
    \draw[->,thick,sourceblue] (tok) -- (qbox) node[midway,above,font=\tiny]{$\mat{W}_Q$};
    \draw[->,thick,targetgreen] (tok) -- (kbox) node[midway,above,font=\tiny]{$\mat{W}_K$};
    \draw[->,thick,quotientorange] (tok) -- (vbox) node[midway,below,font=\tiny]{$\mat{W}_V$};

    % Explanatory labels on the right
    \node[anchor=west,font=\scriptsize,text=sourceblue!90!black] at (7.8,1.35) {Query: ``What information am I searching for?''};
    \node[anchor=west,font=\scriptsize,text=targetgreen!90!black] at (7.8,0.2) {Key: ``What information do I offer to others?''};
    \node[anchor=west,font=\scriptsize,text=quotientorange!90!black] at (7.8,-0.95) {Value: ``What vector content is transferred?''};
  \end{scope}

  % Lane 2: Geometric Dot-Product Scoring and Convex Value Blending
  \begin{scope}[shift={(0,0)}]
    \node[anchor=south,font=\footnotesize\scshape,text=spacegray!90!black] at (6.6,2.3) {2. Dot-Product Alignment and Convex Value Combination};
    \draw[fill=codebg,draw=spacegray!30,rounded corners=4pt] (-0.6,-1.7) rectangle (14.2,2.2);

    % Left: Dot product angle
    \begin{scope}[shift={(1.8,0.2)}]
      \draw[->,gray!50] (-0.3,0) -- (2.4,0);
      \draw[->,gray!50] (0,-0.3) -- (0,1.8);
      \draw[->,sourceblue,line width=1.4pt] (0,0) -- (1.8,0.5) node[right,font=\scriptsize]{$\vec{q}_1$};
      \draw[->,targetgreen,line width=1.4pt] (0,0) -- (1.5,1.4) node[above,font=\scriptsize]{$\vec{k}_2$};
      \draw[gray!70] (0.6,0.17) arc (15:43:0.6);
      \node[font=\tiny] at (0.9,0.45) {$\theta_{12}$};
      \node[font=\scriptsize,align=center] at (1.0,-1.1) {Score: $\vec{q}_1\T\vec{k}_2 / \sqrt{d_k}$\\[3pt]$\propto \cos\theta_{12}$};
    \end{scope}

    % Middle: Softmax normalization arrow
    \draw[->,thick,spacegray!80!black] (4.0,0.2) -- (6.0,0.2)
      node[midway,above,font=\scriptsize,text=structurepurple!90!black]{softmax}
      node[midway,below,font=\tiny]{$\sum \alpha_{1j} = 1$};

    % Right: Convex combination in Value Space
    \begin{scope}[shift={(9.2,0.2)}]
      \draw[quotientorange!60,dashed,line width=1.2pt] (0,1.3) -- (2.4,-0.2);
      \draw[->,quotientorange,line width=1.3pt] (0,0) -- (0,1.3) node[above left,font=\scriptsize]{$\vec{v}_1$};
      \draw[->,quotientorange,line width=1.3pt] (0,0) -- (2.4,-0.2) node[below right,font=\scriptsize]{$\vec{v}_2$};
      \fill[quotientorange] (0,1.3) circle (2pt);
      \fill[quotientorange] (2.4,-0.2) circle (2pt);

      % Resulting blend vector
      \coordinate (blend) at (1.6,0.3); % 0.33 v1 + 0.67 v2
      \draw[->,vecred,line width=1.6pt] (0,0) -- (blend) node[above right,font=\small]{$\text{Out}_1$};
      \fill[vecred] (blend) circle (2.2pt);
      \node[font=\scriptsize,align=center] at (1.2,-1.1) {Convex blend in Value Space:\\$\text{Out}_1 = \alpha_{11}\vec{v}_1 + \alpha_{12}\vec{v}_2$};
    \end{scope}
  \end{scope}
\end{tikzpicture}
\caption{The geometric mechanism of self-attention. Top: Input token vectors are linearly projected into Query, Key, and Value semantic subspaces via learned weight matrices $\mat{W}_Q, \mat{W}_K, \mat{W}_V$. Bottom: The query $\vec{q}_1$ acts as a directional probe, evaluating its dot-product projection against each key $\vec{k}_j$. Softmax normalizes these scores into non-negative weights summing to $1$, which then synthesize a convex combination of value vectors in Value Space.}
\label{ch34:fig-attention}
\end{figure}

\begin{applicationbox}
\textbf{Why transformers scale.} The reason attention powers modern language models is
that at its operational core lies dense matrix multiplication -- $\mat{Q}\mat{K}\T$ and then
$(\cdots)\mat{V}$ -- which is exactly the operation GPUs are built to do astronomically
fast. A model with hundreds of billions of parameters is, in the inner loop, a long
sequence of big matrix multiplies. When you read that a language model was ``trained on
thousands of GPUs,'' the thing those GPUs were multiplying, over and over, were matrices
of queries, keys, values, and weights. The linear algebra in this book is, quite
literally, the computational engine of modern AI.
\end{applicationbox}

\begin{trapbox}
It is tempting to think attention is ``mostly nonlinear and mysterious.'' It is the
opposite: the only nonlinearity in the whole formula is the softmax. Everything else --
forming $\mat{Q},\mat{K},\mat{V}$, scoring with $\mat{Q}\mat{K}\T$, blending with
$\mat{V}$ -- is matrix multiplication. If you can multiply matrices and you understand the
dot product as similarity, you understand the skeleton of a transformer.
\end{trapbox}

\section{Low rank: why a little structure goes a long way}
\label{ch34:lowrank}

We close with the idea that quietly makes large-scale ML feasible: \textbf{low-rank
structure.} It is the SVD chapter cashed out in practice, and it explains
both why models can be compressed and why they generalize at all.

\subsection{Rank as the number of independent directions}

The \textbf{rank} of a matrix is the count of independent directions in its
column space -- the true dimensionality of what it does. A matrix can be huge,
$1000\times1000$, and yet have rank $5$: its million entries are really controlled by only
a handful of underlying patterns. The SVD makes this precise. Any matrix factors as
\[
\mat{A} = \sum_{i=1}^{r} \sigma_i\,\vec{u}_i\vec{v}_i\T,
\]
a sum of rank-one pieces, each a single direction $\vec{u}_i$ scaled by a singular value
$\sigma_i$. The singular values come in decreasing order, $\sigma_1\ge\sigma_2\ge\cdots$,
so the first few terms carry most of the action.

\begin{intuitionbox}
A rank-one matrix $\vec{u}\vec{v}\T$ is the simplest non-trivial matrix there is: every
row is a multiple of $\vec{v}\T$, every column a multiple of $\vec{u}$. It stores
$m+n$ numbers but spreads them across $m\times n$ entries -- pure structure, no
independent detail. The SVD says \emph{every} matrix is a stack of these, and that the
$\sigma_i$ tell you how much each layer matters. ``Low rank'' means a few layers already
tell almost the whole story.
\end{intuitionbox}

\subsection{Low-rank approximation: keep the loud directions}

The payoff is a theorem of Eckart and Young: the \emph{best} rank-$k$
approximation of any matrix -- best in the least-squares sense -- is obtained by keeping
the top $k$ terms of the SVD and throwing the rest away:
\[
\mat{A} \approx \mat{A}_k = \sum_{i=1}^{k} \sigma_i\,\vec{u}_i\vec{v}_i\T.
\]
If the singular values decay quickly (and for real data they almost always do), then a
small $k$ captures almost all of $\mat{A}$ (Figure~\ref{ch34:fig-svdecay}), and you have
compressed an $m\times n$ matrix
($mn$ numbers) down to $k(m+n)$ numbers. For a $1000\times1000$ matrix kept at rank $10$,
that is a 50-fold compression.

\begin{figure}[H]
\centering
\begin{tikzpicture}[>=Stealth]
\begin{axis}[
  width=11cm, height=5cm,
  ybar, bar width=7pt,
  axis lines=left, xlabel={index $i$}, ylabel={singular value $\sigma_i$},
  xmin=0.3, xmax=10.7, ymin=0, ymax=10.5,
  xtick={1,2,3,4,5,6,7,8,9,10}, ytick={2,4,6,8,10},
  enlarge x limits=0.04,
]
  \addplot[fill=sourceblue!55,draw=sourceblue!80!black]
    coordinates {(1,10)(2,6.5)(3,3.2)(4,1.4)(5,0.7)(6,0.4)(7,0.25)(8,0.15)(9,0.1)(10,0.07)};
\end{axis}
\draw[targetgreen,very thick,dashed] (3.03,0) -- (3.03,4.1);
\node[targetgreen!60!black,align=center] at (5.9,3.6)
  {\small keep these $k=3$ ``loud'' directions;\\[-2pt]\small the rest are near-zero noise};
\end{tikzpicture}
\caption{Why low-rank approximation works. For typical data the singular values $\sigma_i$
fall off fast, so a few leading rank-one pieces $\sigma_i\vec{u}_i\vec{v}_i\T$ carry almost
all of the matrix. Truncating after $k$ terms (here $k=3$, dashed line) is the Eckart--Young
best rank-$k$ approximation and discards only the small ``tail.''}
\label{ch34:fig-svdecay}
\end{figure}

\begin{applicationbox}
\textbf{Recommender systems and model compression.} The famous Netflix-style
recommender is a low-rank story: the giant ``users $\times$ movies'' ratings matrix is
modeled as a product of a thin ``users $\times$ taste'' matrix and a thin ``taste
$\times$ movies'' matrix -- a low-rank factorization that fills in the blanks you have
not rated. The same idea now compresses enormous neural networks: a method called LoRA
(``low-rank adaptation'') fine-tunes a giant model by adding only a low-rank update
$\mat{B}\mat{A}$ to its weight matrices, training a few million numbers instead of
billions.
\end{applicationbox}

\subsection{Principal component analysis}

Run the SVD on a centered data matrix and you have \textbf{PCA}, the most-used
dimensionality-reduction method in all of data science. The top singular directions
$\vec{v}_i$ are the \emph{principal components} -- the orthogonal directions along which
your data varies most. Projecting the data onto the first two or three of them is how a
high-dimensional dataset gets drawn as a 2D scatter plot you can actually look at. The
directions of greatest variance are exactly the leading eigenvectors of the covariance
matrix $\mat{X}_c\T\mat{X}_c$ -- the second geometry of your data matrix from
Section~\ref{ch34:data-matrix}, returning to close the loop.

\begin{aestheticbox}
Across this chapter, the dot product gave us \emph{similarity}. Projection gave us
\emph{regression}. Matrix multiplication gave us \emph{layers, attention, and the chain
rule}. The gradient gave us \emph{learning}. The SVD gave us \emph{compression and the
principal directions of data}. Modern machine learning relies on these core linear-algebraic structures, scaled up and run on parallel hardware.
\end{aestheticbox}

\begin{walkbox}
Take a walk and answer these without paper. (1) A neural network with ten layers but no
nonlinearities -- what single, simpler model is it equivalent to, and why?
(2) Cosine similarity ignores vector length; name one situation where that is exactly the
right thing to ignore, and one where it might throw away something you care about.
(3) Gradient descent steps in the direction $-\nabla L$. If the loss surface were a
perfectly round bowl, would the path to the bottom be straight or curved? What about a
long, narrow valley?
\end{walkbox}

\subsection*{Honest limits and computational traps}

At scale, four mathematical bottlenecks confront every machine-learning system.

First, standard scaled dot-product attention carries a quadratic memory wall. To compare
every token to every other token in a sequence of length $N$, the model must materialize
an $N \times N$ matrix of attention scores. For a modern context of $N = 100{,}000$ tokens,
storing that single matrix for every attention head consumes tens of gigabytes of GPU
memory. The search for sub-quadratic, sparse, and low-rank linear attention is an active
race across applied mathematics.

Second, repeated products of matrices are numerically treacherous. Pushing gradients back
through fifty or a hundred layers means evaluating long product chains $\prod_{l=1}^L
\mat{W}_l\T$. As we saw in Chapter 33, if the singular values of these matrices
systematically exceed $1$, gradients explode exponentially into floating-point overflow;
if they fall below $1$, gradients vanish exponentially to machine zero. Modern architectures
survive only by adding identity bypasses (residual skip connections) and normalizing
activations at every layer to keep the spectrum bounded near unity.

Third, the safety of convexity is gone. Unlike linear regression, where the squared error
is a convex quadratic bowl with a guaranteed global minimum reachable via normal equations
or stable QR decomposition, neural loss surfaces twist into non-convex shapes. They are
crowded with saddle points, flat plateaus, and ill-conditioned ravines where gradient
descent can oscillate or stall without second-order curvature information.

Finally, there is the puzzle of internal representations. Although every forward and
backward step is transparent linear algebra, the representations learned in intermediate
layers remain difficult to interpret. High-dimensional vector spaces allow
``superposition,'' where networks pack vastly more concepts into an activation vector than
it has dimensions, aligning distinct features with nearly orthogonal directions. We can
compute the numbers with complete precision, but reading off human-interpretable meaning
from arbitrary coordinates remains an open frontier.

\subsection*{Where we go next in the book}

Throughout this chapter, our data matrix $\mat{X}$ treated each example as an isolated
row in a feature space. In Chapter~35, we connect these points into relational webs,
using adjacency and Laplacian matrices to study networks, random walks, and Markov chains.

\section{Exercises}
\label{ch34:exercises}

\begin{exercise}
A data matrix has rows $(2,4),(4,6),(6,8)$.
(a) Compute the column-mean vector and write the centered matrix $\mat{X}_c$.
(b) Verify each column of $\mat{X}_c$ sums to zero.
(c) Compute the $2\times2$ matrix $\mat{X}_c\T\mat{X}_c$. Are the two features perfectly
correlated? Explain geometrically what you see in the centered points.
\end{exercise}

\begin{exercise}
Fit a line $\hat y = wx+b$ by least squares to the points $(0,1),(1,1),(2,3)$.
(a) Write $\mat{X}$ (with a ones-column) and $\vec{y}$.
(b) Form and solve the normal equations $\mat{X}\T\mat{X}\vec{w}=\mat{X}\T\vec{y}$.
(c) Compute the residual vector and verify it is orthogonal to both columns of $\mat{X}$.
\end{exercise}

\begin{exercise}
Given the embeddings $\vec{a}=\colvec{1\\0\\1}$, $\vec{b}=\colvec{1\\1\\0}$,
$\vec{c}=\colvec{0\\0\\2}$, compute all three pairwise cosine similarities. Rank the pairs
from most to least similar, and say which pair is closest to orthogonal.
\end{exercise}

\begin{exercise}
A single neural-network layer has
$\mat{W}=\begin{bmatrix}2&0\\-1&1\\0&3\end{bmatrix}$, $\vec{b}=\colvec{-3\\0\\-2}$, and
ReLU activation. Compute the output for input $\vec{x}=\colvec{1\\1}$. Which neurons fire?
\end{exercise}

\begin{exercise}
Show that stacking three affine layers $\vec{x}\mapsto\mat{W}_3(\mat{W}_2(\mat{W}_1\vec{x}
+\vec{b}_1)+\vec{b}_2)+\vec{b}_3$ with no nonlinearities equals a \emph{single} affine map
$\mat{W}'\vec{x}+\vec{b}'$. Give explicit formulas for $\mat{W}'$ and $\vec{b}'$. What does
this prove about deep linear networks?
\end{exercise}

\begin{exercise}
(Gradient descent by hand.) Let $L(w)=w^2-4w+5$.
(a) Compute $\nabla L = L'(w)$ and find the minimizer analytically.
(b) Run three steps of gradient descent from $w_0=0$ with learning rate $\eta=0.3$.
(c) Repeat with $\eta=1.2$ and describe what happens. What property of $\eta$ separates
convergence from divergence here?
\end{exercise}

\begin{exercise}
Prove that for a unit vector $\vec{u}$, the directional derivative $\nabla L\cdot\vec{u}$
is maximized when $\vec{u}$ points along $\nabla L$ and minimized when it points along
$-\nabla L$. (Use $\nabla L\cdot\vec{u}=\norm{\nabla L}\cos\theta$.) Explain in one
sentence why this is the entire justification for the gradient-descent update rule.
\end{exercise}

\begin{exercise}
Let $\mat{Q}=\begin{bmatrix}1&0\\0&1\end{bmatrix}$ (two query rows) and
$\mat{K}=\begin{bmatrix}1&0\\1&1\end{bmatrix}$ (two key rows).
(a) Compute the score matrix $\mat{Q}\mat{K}\T$.
(b) Apply softmax to each row (you may leave answers in terms of $e$).
(c) With $\mat{V}=\begin{bmatrix}2&0\\0&2\end{bmatrix}$, compute the attention output.
Interpret which token each query attended to most.
\end{exercise}

\begin{exercise}
A rank-one matrix is $\mat{A}=\vec{u}\vec{v}\T$ with $\vec{u}=\colvec{1\\2\\1}$ and
$\vec{v}=\colvec{2\\1}$.
(a) Write out $\mat{A}$ explicitly (it is $3\times2$).
(b) Verify by inspection that every column is a multiple of $\vec{u}$ and every row a
multiple of $\vec{v}\T$.
(c) How many numbers does it take to store $\mat{A}$ as a rank-one product versus as a
full $3\times2$ table? Generalize to $m\times n$.
\end{exercise}

\begin{exercise}
(Challenge.) The covariance-style matrix $\mat{X}\T\mat{X}$ and the Gram matrix
$\mat{X}\mat{X}\T$ generally have different sizes. Prove that they nonetheless share the
same nonzero eigenvalues. (Hint: if $\mat{X}\T\mat{X}\,\vec{v}=\lambda\vec{v}$ with
$\lambda\neq0$, consider the vector $\mat{X}\vec{v}$ and show it is an eigenvector of
$\mat{X}\mat{X}\T$ with the same $\lambda$.) Explain why this is exactly why PCA can be
computed from whichever of the two matrices is smaller.
\end{exercise}
